\documentclass[10pt,a4paper]{amsart}
\usepackage[margin=19mm]{geometry}
\usepackage{amsmath,amssymb,mathtools}
\usepackage[scr=rsfs,cal=euler]{mathalfa}
\usepackage{xfrac}
\usepackage[unicode,hidelinks,bookmarks=false]{hyperref}
\newcommand{\ie}{i.e.\ }
\newcommand{\cf}{cf.\ }
\newcommand{\resp}{respectively\ }
\newcommand{\Subsection}{\subsection}
\providecommand{\nobreakdash}{\nobreak\hbox{-}}
\setlength{\emergencystretch}{2em}
\multlinegap0pt
\usepackage{bm}
\usepackage{bbm}
\usepackage{dsfont}
\usepackage{xspace}
\usepackage{cancel}
\usepackage{mathtools}
\usepackage{amsthm}
\makeatletter
\@ifundefined{thm}{\newtheorem{thm}{Thm}}{}
\@ifundefined{lemma}{\newtheorem{lemma}[thm]{Lemma}}{}
\makeatother
\title[Travelling Waves in a Mathematical Model for Oncolytic Virot]{Travelling Waves in a Mathematical Model for Oncolytic Virotherapy}
\author{Negar Mohammadnejad, Thomas Hillen}
\date{Source: arXiv:2601.02568v2 (CC BY 4.0)}
\begin{document}
\maketitle

\noindent\emph{Source and licence.} Travelling Waves in a Mathematical Model for Oncolytic Virotherapy. Version arXiv:2601.02568v2 (CC BY 4.0), distributed under the Creative Commons Attribution 4.0 International licence, as recorded in the arXiv licence metadata for this exact version and shown on the abstract page https://arxiv.org/abs/2601.02568v2. Full rights notice: licenses/arxiv-2601.02568-CC-BY-4.0.txt.\par\smallskip

\noindent\textit{Editorial scope.} Counted unit: the Lemma whose source label is \texttt{upper\_condition\_lemma} in the article (source0.tex lines 1091--1099, source label \texttt{upper\_condition\_lemma}), delivered together with the article's own complete original proof of it (source0.tex lines 1101--1363, one \texttt{proof} environment of 263 lines). No mathematics is written, repaired or re-derived: statement and proof are the article's own text line for line. Required context retained uncounted: uppersolutions (lines 1006--1020). Every definition, display and label of those lines is retained unchanged. Omitted from this package and not redistributed: 1--1005, 1021--1090, 1100--1100, 1364--2639. Nothing inside a retained excerpt is omitted or altered: every retained body is delivered exactly as the article's lines are written, so no omission receipt of any kind accompanies this package. Citations: the retained excerpts carry no citation command at all, so no bibliography entry of the article is transcribed and no reference is redistributed. Reference adaptation: none was needed and none was made; every reference inside the delivered unit resolves inside it, so no unresolved reference or citation is produced. Printed slips kept literal: no printed word, symbol, hypothesis or number of the counted statement or of its proof was altered. Layout and class: the article's own class is \texttt{article} with packages including graphicx, fontenc, inputenc, lmodern, xcolor, placeins, multirow, amsmath, amsfonts, amsthm, amssymb, comment, pgfplots, color; the wrapper is the lane's pinned \texttt{amsart} editorial wrapper at 10pt on A4, which restates the theorem-like environments the retained text needs and adds the article's own macro definitions that the retained text needs (none), with extra wrapper packages (bm, bbm, dsfont, xspace, cancel, mathtools). The delivered numbering is the unit's own. Assets: no retained span contains a figure environment, a graphics-inclusion command, a PDF-page inclusion, a table or a TikZ picture; no image file is copied into this package, the package declares no asset dependency and no image file is redistributed with it. Licence: version arXiv:2601.02568v2 of this article is distributed under the Creative Commons Attribution 4.0 International licence, as recorded in the arXiv licence metadata for this exact version and shown on the abstract page https://arxiv.org/abs/2601.02568v2. A full rights notice is delivered at licenses/arxiv-2601.02568-CC-BY-4.0.txt. Characters: every retained excerpt is pure ASCII, so the delivered engine, XeLaTeX with its default Latin Modern text font, reports no missing character.
Thus, we can rewrite the constructed upper solutions as \begin{align}\label{uppersolutions}
     \nonumber \bar{B}(\xi)& =
    \begin{cases}
     e^{\rho \xi} \dfrac{(c \rho+a+1-D \rho^2)}{(c \rho +a- D \rho^2)(c\rho +1- D\rho^2)},& \quad \xi \leq \bar{\xi_1}  \\ 
        1, & \quad \xi > \bar{\xi}_1
    \end{cases} \\\bar{I}(\xi) &=
    \begin{cases}
       e^{\rho \xi}\dfrac{1}{c \rho +a - D \rho^2}, & \quad \xi \leq \bar{\xi_2} \\ 
        1, & \quad \xi > \bar{\xi_2}
    \end{cases}\\ \nonumber\bar{V}(\xi)& =
    \begin{cases}
        e^{\rho \xi}, & \quad \xi \leq\bar{\xi_3} \\ 
         \dfrac{\theta}{\gamma},\, & \quad \xi >\bar{\xi_3},
    \end{cases}
\end{align}

\begin{lemma}
    For $c> \bar{c}$ and $\bar{B}, \bar{I}$ and $\bar{V}$ as given above \eqref{uppersolutions}, the following inequalities are satisfied. 
    \begin{align}\label{upper_condition_lemma}
      \nonumber  c \bar{B}'- D \bar{B}'' 
 &\geq- \bar{B}+\bar I +\bar{B}^2-\bar{B} \bar{I}+ \bar{V}-\bar{B} \bar{V}, \\  c \bar{I}'- D \bar{I}'' 
 &\geq -a\bar{I}+ (1-\bar{B})\bar{V}, \\ \nonumber c \bar{V}'-  \bar{V}'' 
 &\geq \theta \bar{I} - \gamma \bar{V}.
    \end{align}
\end{lemma}

\begin{proof} 
\textbf{Step 1.} We first show that under the condition $c > D\rho$, the upper solution 
$\bar{U} = (\bar{B}, \bar{I}, \bar{V})$ satisfies the linearization inequality
\begin{align}\label{linearization_inequality}
    cU' - \mathcal{D}U'' - f(U) \;\ge\; cU' - \mathcal{D}U'' - Uf'(0).
\end{align}

For the components $\bar{I}$ and $\bar{V}$ the inequality follows directly.  
For $\bar{I}$ we compute
\begin{align}
    c\bar{I}' - D\bar{I}'' - f(\bar{I})
    &= c\bar{I}' - D\bar{I}'' + a\bar{I} + (\bar{B}-1)\bar{V} \\
    &\ge c\bar{I}' - D\bar{I}'' + a\bar{I} - \bar{V}
      = c\bar{I}' - D\bar{I}'' - \bar{I}f'(0),
\end{align}
since $\bar{B} \ge 0$.  
For $\bar{V}$ we similarly obtain
\begin{align}
    c\bar{V}' - \bar{V}'' - f(\bar{V})
    = c\bar{V}' - \bar{V}'' + \gamma \bar{V} - \theta \bar{I}
    = c\bar{V}' - \bar{V}'' - \bar{V}f'(0).
\end{align}

For the component $\bar{B}$, we have
\begin{align}
    c\bar{B}' - D\bar{B}'' - f(\bar{B})
    &= c\bar{B}' - D\bar{B}'' 
    + \bar{B} - \bar{I} - \bar{B}^2 + \bar{B}\bar{I} - \bar{V} + \bar{B}\bar{V} \\
    &= c\bar{B}' - D\bar{B}'' 
       + \bar{B}\bigl(-\bar{B} + \bar{I} + \bar{V}\bigr)
       + (\bar{B}-\bar{I} - \bar{V}) \\
    &\ge c\bar{B}' - D\bar{B}'' + \bar{B} - \bar{I} - \bar{V}
      = c\bar{B}' - D\bar{B}'' - \bar{B}f'(0),
\end{align}
which holds provided that
\begin{align} \label{I+V>B}
    \bar{B} \le \bar{I} + \bar{V}.
\end{align}

For $\xi < \bar{\xi}_1$, substituting the expressions from~\eqref{uppersolutions} into~\eqref{I+V>B} gives
\begin{align}
    \bar{I} + \bar{V}
    = \frac{c\rho^{2}+a+1 - D\rho^2}{c\rho+a - D\rho^2}
    \;\ge\;
    \frac{(c\rho+a+1 - D\rho^2)}
         {(c\rho+a - D\rho^2)(c\rho+1 - D\rho^2)}=\bar{B}.
\end{align}

Since we have
\[
c\rho+a - D\rho > 0,
\qquad
c\rho+1 - D\rho^2 > 0,
\]
the inequality reduces to the requirement
\begin{align}\label{upper_requirement}
    0 \le c\rho - D\rho^{2},
\end{align}
which is equivalent to the condition $c \ge D\rho$ for $\rho>0$.

On the interval $\bar{\xi}_1 < \xi < \bar{\xi}_2$, only the condition at 
$\xi = \bar{\xi}_1$ needs to be verified, because $\bar{B}$ is identically 
equal to $1$ on this interval, while both $\bar{I}$ and $\bar{V}$ are 
increasing functions.  
At $\xi = \bar{\xi}_1$ we have

\begin{align}
   \bar{B}= e^{\rho \bar{\xi}_1}
    \left(
     \frac{(c\rho+a+1 - D\rho^2)}
         {(c\rho+a - D\rho^2)(c\rho+1 - D\rho^2)}
    \right)
    = 1,
\end{align}
which simplifies to
\begin{align}
    e^{\rho \bar{\xi}_1}
    \left(
     \frac{(c\rho+a+1 - D\rho^2)}
         {(c\rho+a - D\rho^2)}
    \right)
    = c\rho - D\rho^{2} + 1 \;\ge\; 1,
\end{align}
since $c > D\rho$.  
Thus,
\begin{align}
    \bar{I} + \bar{V}
    = e^{\rho \bar{\xi}_1}
    \left(
        \frac{(c\rho+a+1 - D\rho^2)}
         {(c\rho+a - D\rho^2)}
    \right)
    \ge 1.
\end{align}

For all $\xi > \bar{\xi}_2$, the inequality $\bar{I} + \bar{V} \ge \bar{B}$ holds trivially because  
$\bar{I} = 1$ and $\bar{V} \ge 0$.

Since all components $\bar{B},\ \bar{I},$ and $  \bar{V}$ satisfy the linearization inequality $cU' - \mathcal{D}U'' - f(U) \;\ge\; cU' - \mathcal{D}U'' - Uf'(0)$, instead of proving the inequalities \eqref{upper_condition_lemma}, we can equivalently prove
\begin{align}
    c \bar{B}' - D\bar{B}'' &\ge -\bar{B} + \bar{I} + \bar{V}, \\
    c \bar{I}' - D\bar{I}'' &\ge -a\bar{I} + \bar{V}, \\
    c \bar{V}' - \bar{V}'' &\ge \theta \bar{I} - \gamma \bar{V}.
\end{align}
Thus, in more complicated cases, it is sufficient to verify the linearized inequality
\[
cU' - \mathcal{D}U'' - Uf'(0) \;\ge\; 0,
\]
instead of the full nonlinear inequality
\[
cU' - \mathcal{D}U'' - f(U) \;\ge\; 0.
\]
In the following steps, we prove the inequalities \eqref{upper_condition_lemma} for all possible positions of $\xi$ within the ordering \[\bar{\xi}_1\le \bar{\xi}_2 \le \bar{\xi}_3.\]
\textbf{Step 2.} We first consider the region $\xi < \bar{\xi}_1$, where the upper solutions are given by 
\[
\bar{B}(\xi)= e^{\rho \xi}\,
\frac{(c\rho+a+1 - D\rho^2)}
         {(c\rho+a - D\rho^2)(c\rho+1 - D\rho^2)}, 
\qquad
\bar{I}(\xi)= e^{\rho \xi}\,\frac{1}{c\rho+a-D\rho^2}, 
\qquad
\bar{V}(\xi)= e^{\rho \xi}.
\]

Substituting these into the linearized equation
\[
cU' - D U'' - U f'(0),
\]
we compute each component. For $\bar{B}$, we obtain
\begin{align}
c\bar{B}' - D\bar{B}'' + \bar{B} - \bar{I} - \bar{V}
= e^{\rho \xi}\left[
\frac{(c\rho +1- D\rho^{2})(c\rho+a+1-D\rho^2)}{(c\rho+a-D\rho^2)(c\rho+1-D\rho^2)}
-\frac{1}{c\rho+a-D\rho^2} - 1
\right]
=0.
\end{align}

For $\bar{I}$, we have
\begin{align}
c\bar{I}' - D\bar{I}'' + a\bar{I} - \bar{V}
= e^{\rho \xi}\left[
\frac{D\rho^{2}-a-c\rho}{D\rho^{2}-a-c\rho} - 1
\right]
=0.
\end{align}

For $\bar{V}$, substitution yields
\begin{align}
c\bar{V}' - \bar{V}'' - \theta \bar{I} + \gamma\bar{V}
= e^{\rho \xi}\left[
(c\rho - \rho^{2} + \gamma)(D\rho^{2}-a-c\rho) + \theta
\right]
=0.
\end{align}

The last equality follows from the relation  
\begin{align}\label{det_identitty_lemma_upper_proof}
0 = \det(A(\rho) - c\rho\,\mathbb{I})
     = (D\rho^{2}-1)\left[(c\rho - \rho^{2} - \gamma)(D\rho^{2}-a-c\rho) + \theta\right],
\end{align}
since $c\rho$ is the eigenvalue of $A(\rho)$.


\textbf{Step 3.} Next, consider $\bar{\xi}_1 \le \xi < \bar{\xi}_2 < \bar{\xi}_3$, where  
\[
\bar{B}=1, 
\qquad 
\bar{I}= e^{\rho \xi}\frac{1}{c\rho+a-D\rho^2},
\qquad 
\bar{V}= e^{\rho \xi}.
\]

For $\bar{B}$, we have
\begin{align}
c\bar{B}'-D\bar{B}''+\bar{B}-\bar{B}^{2}+\bar{B}\bar{I}+\bar{B}\bar{V}-\bar{I}-\bar{V}
= 1 - 1 + \bar{I} + \bar{V} - \bar{I} - \bar{V}
=0.
\end{align}

For $\bar{I}$, we obtain
\begin{align}
c\bar{I}' - D\bar{I}'' + a\bar{I} - \bar{V}
= e^{\rho \xi}\left[
(c\rho+a - D\rho^{2} )\frac{1}{c\rho+a-D\rho^2} - 1
\right]
=0.
\end{align}

For $\bar{V}$, we compute
\begin{align}
c\bar{V}' - \bar{V}'' - \theta \bar{I} + \gamma\bar{V}
= e^{\rho \xi}\left[
(c\rho - 1 + \gamma) + \frac{\theta}{D\rho^{2}-a-c\rho}
\right]
=0.
\end{align}

Once again, the final equality results from the identity $\det(A(\rho)-c\rho \mathbb{I})=0$.

\textbf{Step 4.} In the next interval, $\bar{\xi}_1 < \bar{\xi}_2 \le \xi < \bar{\xi}_3$, the upper solutions are  
\[
\bar{B}=1, \qquad
\bar{I}= 1,
\qquad 
\bar{V}=e^{\rho \xi}.
\]

For $\bar{B}$ we again have
\begin{align}
c\bar{B}'-D\bar{B}''+\bar{B}-\bar{B}^{2}+\bar{B}\bar{I}+\bar{B}\bar{V}-\bar{I}-\bar{V}
= 1 - 1 + \bar{I} + \bar{V} - \bar{I} - \bar{V}
=0.
\end{align}

For $\bar{I}$,
\begin{align}
c\bar{I}' - D\bar{I}'' + a\bar{I} -(1-\bar{B}) \bar{V}
= a
\ge 0,
\end{align}
and this inequality holds since $a\ge0.$

For $\bar{V}$, noting that \eqref{det_identitty_lemma_upper_proof} gives $(c\rho - \rho^{2} - \gamma)=\dfrac{\theta}{(c \rho+a-D\rho^{2})}$, we write
\begin{align}
\bar{V}'-D\bar{V}''- \theta \bar{I} + \gamma \bar{V}
=& e^{\rho \xi}\left(c \rho-\rho^2+\gamma \right) - \theta
\\=& e^{\rho \xi}\left(\dfrac{\theta}{c\rho+a-D\rho^2}\right) - \theta =\theta\left(e^{\rho \xi}\dfrac{1}{c \rho+a-D\rho^2}-1\right)
\ge 0,
\end{align}
because in this region $\bar{I}=1$ which means $1\le e^{\rho \xi} \dfrac{1}{c \rho+a-D\rho^2}$.


\textbf{Step 5.} Finally, consider the last interval, $\xi \ge \bar{\xi}_3$, where  
\[
\bar{B}=1, 
\qquad 
\bar{I}=1,
\qquad 
\bar{V}=\frac{\theta}{\gamma}.
\]

For $\bar{B}$, we again have
\begin{align}
c\bar{B}'-D\bar{B}''+\bar{B}-\bar{B}^{2}+\bar{B}\bar{I}+\bar{B}\bar{V}-\bar{I}-\bar{V}
=0.
\end{align}

For $\bar{V}$,
\begin{align}
\bar{V}'-D\bar{V}''- \theta \bar{I} + \gamma \bar{V}
= \gamma\left(\frac{\theta}{\gamma}\right) - \theta
\ge0.
\end{align}

For $\bar{I}$,
\begin{align}
c\bar{I}'-D\bar{I}'' + a\bar{I} - (1-\bar{B})\bar{V}
= a - 0
= a \ge 0.
\end{align}

\end{proof}

\end{document}
